\appendix
\chapter{Appendices}
\label{chap:appendix}

\section{Anomaly template CSV format}
\label{app:csv-format}

Each anomaly is encoded as a CSV file in
\texttt{data-gen/anomalies/}. The columns are described below.

\begin{table}[H]
    \centering
    \footnotesize
    \begin{tabular}{lp{0.62\linewidth}}
        \toprule
        Column & Meaning \\
        \midrule
        \texttt{time}    & Normalised time of the keypoint, in $[0, 1]$. \\
        \texttt{value}   & Normalised amplitude of the keypoint, in $[0, 1]$. \\
        \texttt{t\_minus}, \texttt{t\_plus} & Boolean flags allowing
        the keypoint to jitter in the negative or positive time direction. \\
        \texttt{v\_minus}, \texttt{v\_plus} & Same, for the amplitude axis. \\
        \texttt{interp}  & Interpolation mode for the segment that
        \emph{starts} at this keypoint:
        \texttt{linear} (straight line),
        \texttt{exp\_up} (concave-up exponential),
        \texttt{exp\_down} (concave-down exponential),
        \texttt{bell} (smoothed half-cosine). \\
        \bottomrule
    \end{tabular}
    \caption{Schema of the anomaly-template CSV files.}
    \label{tab:csv-schema}
\end{table}

\paragraph{Example: \texttt{bell.csv}.}
\begin{lstlisting}[caption={Encoding of the \texttt{bell} template.}, label={lst:bell-csv}]
time,value,t_minus,t_plus,v_minus,v_plus,interp
0.0,0.0,0,0,0,0,bell
0.5,1.0,1,1,1,1,bell
1.0,0.0,0,0,0,0,linear
\end{lstlisting}

The first and last keypoints have all jitter flags off, so the
template always starts and ends at $(0,0)$ and $(1,0)$. The middle
keypoint can jitter in both time and amplitude; the two half-cosine
segments make the curve smooth.

\section{Hyperparameter reference}
\label{app:hp}

\subsection{Data generator}
\label{app:hp-gen}

\Cref{tab:v1-vs-v2-config} compares the v1 and v2 generator
configurations side by side.

\begin{table}[H]
    \centering
    \footnotesize
    \begin{tabular}{lcc}
        \toprule
        Parameter & v1 & v2 \\
        \midrule
        Total points          & $10^{7}$ & $10^{7}$ \\
        Chunk size            & $5 \cdot 10^{4}$ & $5 \cdot 10^{4}$ \\
        \midrule
        Number of injections  & $\sim [5{,}000, 8{,}000]$
                              & $\sim [1{,}300, 1{,}700]$ \\
        Injection jitter      & $\pm 200$ samples
                              & $\pm 400$ samples \\
        Anomaly density       & $\sim 7$ / $10$k samples
                              & $\sim 1.5$ / $10$k samples \\
        \midrule
        Amplitude range       & uniform in $\sigma \cdot [2.5, 7.6]$
                              & log-uniform in $[15, 40]$ counts (abs.) \\
        Duration range        & uniform in $[200, 800]$ samples
                              & log-uniform in $[30, 1000]$ samples \\
        Deformation $\nu$     & uniform in $[0.02, 0.10]$
                              & uniform in $[0.02, 0.10]$ \\
        \midrule
        HF noise model        & IID $\mathcal{N}(0, \sigma^2)$
                              & AR(1) $\varphi = 0.86$,
                                $\sigma_\varepsilon = 1.12$ \\
        HF noise std $\sigma$ & $2.53$ & $2.30$ (marginal of AR(1)) \\
        OU $\theta$           & $4 \cdot 10^{-6}$
                              & $2.5 \cdot 10^{-5}$ \\
        OU $\sigma_{\mathrm{OU}}$ & $0.0373$
                              & $0.007$ \\
        Global mean $\mu$     & $99.26$ (fixed)
                              & $99.3 \pm 1.5$ (per-recording) \\
        \midrule
        Templates             & 6 (\texttt{bell}, \texttt{bell\_sq},
                                \texttt{fast\_ascend},
                                \texttt{fast\_descend},
                                \texttt{M}, \texttt{square})
                              & 3 (\texttt{bell}, \texttt{fast\_ascend},
                                \texttt{fast\_descend}) \\
        Classes $K$           & 7 & 4 \\
        \bottomrule
    \end{tabular}
    \caption{Configuration of the v1 and v2 generators.}
    \label{tab:v1-vs-v2-config}
\end{table}

\subsection{Training}
\label{app:hp-train}

\begin{table}[H]
    \centering
    \footnotesize
    \begin{tabular}{ll}
        \toprule
        Hyperparameter & Value \\
        \midrule
        Window size $W$                          & $512$ samples \\
        Sliding-window stride at training        & $32$ samples \\
        Sliding-window stride at inference       & $10$ samples \\
        Window labelling threshold $\tau_{\text{cov}}$ & $30\%$ of an event's samples captured \\
        Per-window normalisation                 & subtract window mean, divide by global $\sigma_{\mathrm{train}}$ \\
        $\sigma_{\mathrm{train}}$ (v1)           & $4.21$ \\
        $\sigma_{\mathrm{train}}$ (v2)           & $2.88$ \\
        Train/val/test split                     & temporal $80/20$, then random $80/20$ on train \\
        Optimizer                                & Adam, $\beta_1 = 0.9$, $\beta_2 = 0.999$, $\varepsilon = 10^{-8}$ \\
        Initial learning rate                    & $10^{-3}$ \\
        Batch size                               & $128$ \\
        Maximum epochs                           & $35$ \\
        LR scheduler                             & \texttt{ReduceLROnPlateau}(patience=$2$, factor=$0.5$) \\
        Early stopping                           & patience $10$, min delta $5 \cdot 10^{-3}$ \\
        Loss                                     & weighted categorical cross-entropy \\
        Class weights (v1)                       & $(0.46, 1.18, 1.27, 1.32, 1.24, 1.15, 1.31)$ \\
        Class weights (v2)                       & $(0.27, 8.51, 9.37, 9.63)$ \\
        \bottomrule
    \end{tabular}
    \caption{Shared training configuration.}
    \label{tab:train-config}
\end{table}

\subsection{Inference and event extraction}
\label{app:hp-inf}

\begin{table}[H]
    \centering
    \footnotesize
    \begin{tabular}{ll}
        \toprule
        Hyperparameter & Value \\
        \midrule
        Sliding-window stride & $10$ samples \\
        Inference batch size  & $256$ windows \\
        Smoothing window      & $30$-sample rolling mean on the per-sample probability trace \\
        Background-probability threshold for an event & $\hat{p}_0 < 0.3$ \\
        Tolerance gap between adjacent events & $50$ samples \\
        \bottomrule
    \end{tabular}
    \caption{Configuration of the inference engine and event extractor.}
    \label{tab:inf-config}
\end{table}

\section{Reproducing the figures in this report}
\label{app:reproduce}

The repository layout is:
\begin{lstlisting}[basicstyle=\ttfamily\footnotesize, caption={Top-level project structure.}]
pole-anomaly/
  biblio/                  - the project brief and reference papers
  data-gen/                - v2 data generation and noise reanalysis
    anomaly_generator.py     - template deformation + interpolation
    generate_custom_dataset.py - v2 generator (AR(1) noise, 3 templates)
    analyze_noise_model.py   - AR(p), Ljung-Box, drift, stability
    labeling_view.py         - hand-labelling helper
    visualize_anomalies_templates.py
    test_generator_deformation.py
    anomalies/               - 3 v2 template CSV files
    data/                    - real recordings + hand labels + v2 synth CSV
    logs/                    - PNG/HTML inspection plots
  architectures/           - models, training, inference (v2-trained)
    train.py                 - shared training pipeline (selects model)
    visualize_inference.py   - synthetic-test event-level evaluation
    inference_real_data.py   - real-data deployment trial
    robustness_test.py       - noise/deformation sweep
    inference_engine.py      - sliding-window + event matching
    model_selection.py       - model registry & CLI selector
    saved_models/            - 1d_cnn.py, resNet.py, cnn_attention.py
    logs/                    - per-model training/inference outputs
  pole-anomaly-old/        - v1 generator + v1-trained models
    data-gen/                - 6 templates, IID Gaussian noise generator
    architectures/           - v1-trained model checkpoints + logs
  documentation/           - this report
    main.tex, preamble.tex, sections/*.tex
    figures/build_figures.py - reproduces all PNGs in this report
\end{lstlisting}

To rebuild every figure in this report from scratch:

\begin{lstlisting}[basicstyle=\ttfamily\footnotesize]
# Generate v2 synthetic dataset (~10 minutes; 600 MB)
py data-gen/generate_custom_dataset.py

# (Optional: re-run noise analysis)
py data-gen/analyze_noise_model.py

# (Optional: re-train v2 models)
py architectures/train.py --model 1d_cnn
py architectures/train.py --model resnet
py architectures/train.py --model cnn_attention

# (Optional: re-run synthetic-test and real-data inference)
py architectures/visualize_inference.py --model 1d_cnn
py architectures/inference_real_data.py --model 1d_cnn
py architectures/robustness_test.py --model 1d_cnn
# ... same for resnet, cnn_attention

# Rebuild the report's figures (uses both new and old logs)
py documentation/figures/build_figures.py

# Compile the report
cd documentation && pdflatex main && pdflatex main
\end{lstlisting}

\section{Software environment}
\label{app:env}

All experiments were run on a single laptop with the following
software stack:
\begin{itemize}[leftmargin=*, itemsep=0.2em]
    \item Windows 11, Python $3.13$.
    \item PyTorch $2.12$ (CUDA 12.8), running on an NVIDIA RTX 5060
    Laptop GPU.
    \item NumPy, Pandas, Matplotlib, SciPy, scikit-learn, statsmodels
    for data handling, statistics and reporting.
    \item Plotly for the interactive HTML diagnostics shipped under
    \texttt{logs/}.
\end{itemize}
A pinned dependency list is included in the top-level
\texttt{requirements.txt} of the repository.
